\section{Inversão de Möbius}

\textbf{Problema base.} Inverter somas sobre divisores e contar pares com condição de $\gcd$ (coprimos, $\gcd = x$).

\textbf{Fórmula geral.}
\[
\boxed{g(n) = \sum_{d \mid n} f(d) \iff f(n) = \sum_{d \mid n} \mu(d)\, g\!\left(\tfrac{n}{d}\right)}
\]
onde $\mu$ é a função de Möbius:
\begin{itemize}
  \item $\mu(1) = 1$;
  \item $\mu(d) = 0$ se algum primo aparece ao quadrado em $d$ ($p^2 \mid d$);
  \item $\mu(d) = (-1)^k$ se $d$ é produto de $k$ primos distintos.
\end{itemize}

\textbf{Propriedade fundamental.} Com $[P] = 1$ se $P$ é verdadeira e $0$ se não, ela troca a condição ``coprimos'' por
uma soma sobre os divisores comuns:
\[
[\gcd(a, b) = 1] = \sum_{d \mid \gcd(a, b)} \mu(d) = \sum_{d \mid a,\ d \mid b} \mu(d)
\]

\textbf{Roteiro.} Escreva a contagem como uma soma de $[\gcd = 1]$, troque pela soma de $\mu(d)$ e inverta a ordem das
somas: para cada $d$, só sobra contar quantos objetos têm tudo múltiplo de $d$.

\subsection*{Pares com $\gcd = x$ num vetor}

Seja $C(v)$ a quantidade de elementos múltiplos de $v$ (crivo em $O(M \log M)$, com $M$ o maior valor).
$\gcd(A_i, A_j) = x$ quando os dois são múltiplos de $x$ e, divididos por $x$, ficam coprimos:
\[ \text{Pares} = \sum_{\substack{i, j \\ x \mid A_i,\ x \mid A_j}} \left[\gcd\!\left(\tfrac{A_i}{x}, \tfrac{A_j}{x}\right) = 1\right]. \]
Pela propriedade, com $k$ dividindo $\frac{A_i}{x}$ e $\frac{A_j}{x}$, ou seja, $kx$ dividindo $A_i$ e $A_j$:
\[ \text{Pares} = \sum_{\substack{i, j \\ x \mid A_i,\ x \mid A_j}} \ \sum_{kx \mid A_i,\ kx \mid A_j} \mu(k). \]
Invertendo a ordem, cada $k$ soma $\mu(k)$ uma vez para cada par com os dois múltiplos de $kx$, e há $C(kx)^2$ desses pares:
\[ \text{Pares} = \sum_{k=1}^{\lfloor M/x \rfloor} \mu(k)\, C(kx)^2 \qquad (O(M/x) \text{ por consulta}). \]
São pares ordenados, contando $i = j$; para pares $i < j$, troque $C^2$ por $\binom{C}{2}$.

\subsection*{Pares em $[1, n] \times [1, m]$ com $\gcd = x$}

Os pares com $\gcd = x$ são $(x i', x j')$ com $i', j'$ coprimos, $i' \le n/x$ e $j' \le m/x$:
\[ \text{Pares} = \sum_{i \le n,\ j \le m} [\gcd(i, j) = x] = \sum_{i' \le n/x} \ \sum_{j' \le m/x} [\gcd(i', j') = 1]. \]
Pela propriedade e invertendo a ordem, cada $d$ conta os pares com $d \mid i'$ e $d \mid j'$: há
$\lfloor \frac{n}{xd} \rfloor$ múltiplos de $d$ até $\frac{n}{x}$ e $\lfloor \frac{m}{xd} \rfloor$ até $\frac{m}{x}$:
\[ \text{Pares} = \sum_{d \ge 1} \mu(d) \left\lfloor \tfrac{n}{xd} \right\rfloor \left\lfloor \tfrac{m}{xd} \right\rfloor. \]
Com $x = 1$ são os pares coprimos. Roda em $O(\sqrt{n} + \sqrt{m})$ agrupando os $d$ com os mesmos quocientes e usando prefix sum de $\mu$.

\subsection*{Subconjuntos com $\gcd = 1$ num vetor}

Com o mesmo $C$ e somando sobre os subconjuntos $S$ não vazios do vetor:
\[ \text{Subconjuntos} = \sum_{S} [\gcd(S) = 1] = \sum_{S} \ \sum_{d \mid \gcd(S)} \mu(d). \]
$d \mid \gcd(S)$ quando $d$ divide todos os elementos de $S$. Invertendo a ordem, cada $d$ conta os subconjuntos não vazios
formados só por múltiplos de $d$, que são $2^{C(d)} - 1$:
\[ \text{Subconjuntos} = \sum_{d = 1}^{M} \mu(d) \left(2^{C(d)} - 1\right). \]
